Le système considéré est la caisse ; son mouvement est étudié dans le référentiel terrestre.
Les deux positions considérées sont la position initiale $A\left(v_A=0\right)$ et la position $B$ atteinte par la caisse après le déplacement $\ell=5,00 \mathrm{~m}$ ( $v_B$ à calculer).
La variation d'énergie cinétique entre $A$ et $B$ vaut :

$$
\Delta E_c=\frac{1}{2} m v_B^2-\frac{1}{2} m v_A^2=\frac{1}{2} m v_B^2 .
$$


Les forces appliquées à la caisse sont :
- son poids $\vec{P}_{\text {vertical ; }}$
- la réaction du plan incliné $\vec{R}_N$ qui est normale au support car le contact est sans frottement.
Travail des forces :
- $W_{A B}\left(\vec{R}_N\right)=0$ car $\vec{R}_N$ est constamment perpendiculaire au déplacement.
- $W_{A B}(\vec{P})=+P h=+m g h$.

Le calcul de $b$ s'effectue dans le triangle coloré en bleu :

$$
\sin \alpha=\frac{b}{\ell} \Rightarrow b=\ell \sin \alpha .
$$


Application de la relation :

$$
\begin{gathered}
\Delta E_c=\Sigma W_{A B}\left(\vec{F}_{\text {ext }}\right) \Rightarrow \frac{1}{2} m v_B^2=m g \ell \sin \alpha \\
v_B=\sqrt{2 g \ell \sin \alpha}
\end{gathered}
$$


Numériquement : $v_B=\sqrt{2 \times 9,8 \times 5,0 \times \sin 15^{\circ}} \approx \mathbf{5 , 0} \mathbf{m} \cdot \mathbf{s}^{\mathbf{- 1}}$.
